% ============================================================
% 问题1 小节（国赛风格重写版）—— 直接替换 001sec.txt 中
% 从 \subsection{问题1：...} 到 \subsection{问题2：...} 之前的整块内容。
% 所有数值、公式、图表与引用与 A1-codex 已验证结果一致，未作任何改动。
% ============================================================
\subsection{问题1：预热平衡阶段药材温度、水分浓度变化模型}

\subsubsection{问题分析}

针对问题1，需建立预热平衡阶段（$0\sim1800\,\mathrm{s}$）药材温度场与水分浓度场的
非稳态变化模型，求解不同径向位置处的温度 $T(r,t)$ 与水分浓度 $C(r,t)$，并按题面
表 1、表 2 的格式给出 7 个时刻（$100\sim1800\,\mathrm{s}$）、5 个径向位置
（$0\sim2\,\mathrm{cm}$）处的结果；同时将 $1800\,\mathrm{s}$ 内每隔 $1\,\mathrm{s}$、
每隔 $0.1\,\mathrm{cm}$ 的完整结果保留 4 位小数写入 result1.xlsx（模板见附件 3）。

由题面可知，药材为长 $25\,\mathrm{cm}$、半径 $2\,\mathrm{cm}$ 的细长圆柱
（长径比 $6.25$），侧面受热风均匀吹拂，且题目要求的输出自变量为“到药材中心的
距离”，故本文将问题简化为\textbf{轴对称一维径向热传导--水分扩散模型}。按附录 2，
问题 1 的密度、比热容与导热系数均为常数，水分扩散系数仅随水分浓度变化
（不含温度项），因此温度场与水分场仅通过表面边界条件耦合，两个方程可分别求解。

为判断考察窗内两场的建立程度，求解前先进行特征时间与 Biot 数量级分析。热扩散率与
初始扩散系数分别为 $\alpha=k/(\rho c_p)=1.6886\times10^{-7}\ \mathrm{m^2/s}$ 与
$D(C_0)=7\times10^{-9}e^{-0.89/2.55}=4.9377\times10^{-9}\ \mathrm{m^2/s}$，得特征时间
$\tau_T=R^2/\alpha=2368.9\ \mathrm{s}\approx39.5\ \mathrm{min}$、
$\tau_C=R^2/D(C_0)=8.10\times10^{4}\ \mathrm{s}\approx22.5\ \mathrm{h}$，以及
$Bi=hR/k=1.389$、$Bi_m=h_mR/D(C_0)=3.240$。特征时间与 Biot 数是刻画温度场、水分场
时间尺度与表面控制机制的常用判据\cite{yan2026heat,srikiatden2007moisture}。据此可知：
$\tau_C/\tau_T\approx34$，30 min 约为 $0.76\,\tau_T$ 而仅占 $\tau_C$ 的 $2.2\%$，
故温度场在考察窗内显著建立，水分仅在表面薄层内变化；$Bi=O(1)$，集总参数法不适用，
径向温差可达数摄氏度；$Bi_m=3.24>1$，失水由对流传质边界驱动，将形成表面干、内部湿
的分布。上述量级分析为模型建立与结果解读提供了基本依据。

\subsubsection{模型建立}

预热平衡阶段药材与烘房热空气进行非稳态热--质交换：热量经表面对流换热进入药材并沿
径向传导，水分在浓度梯度作用下向表面扩散、经对流传质进入空气
\cite{yan2026heat,srikiatden2007moisture}。在轴对称一维径向假设下，温度场与水分浓度场
的控制方程为
\begin{equation}
  \rho c_p\frac{\partial T}{\partial t}
  =\frac{1}{r}\frac{\partial}{\partial r}\Big(k r\frac{\partial T}{\partial r}\Big),
  \qquad 0<r<R,\ t>0,
  \label{eq:heat}
\end{equation}
\begin{equation}
  \frac{\partial C}{\partial t}
  =\frac{1}{r}\frac{\partial}{\partial r}\Big(D(C)\,r\frac{\partial C}{\partial r}\Big),
  \qquad 0<r<R,\ t>0,
  \label{eq:moist}
\end{equation}
其中水分扩散系数按附录 2 取经验公式 $D(C)=7\times10^{-9}\exp(-0.89/C)$（指数为
$-0.89/C$）。初始条件为 $T(r,0)=28\,^{\circ}\mathrm{C}$、
$C(r,0)=2.55\ \mathrm{kg/kg}$；边界条件为圆心处
$\partial T/\partial r(0,t)=\partial C/\partial r(0,t)=0$，表面为第三类对流换热与
对流传质：
\begin{equation}
  -k\frac{\partial T}{\partial r}(R,t)=h\big[T(R,t)-T_a(t)\big],
  \qquad
  -D(C_s)\frac{\partial C}{\partial r}(R,t)=h_m\big[C_s(t)-C_a(t)\big],
  \label{eq:bcR}
\end{equation}
式中 $T_a(t)$、$C_a(t)$ 为附件 1 的烘房温度与水分浓度时变序列，$C_s(t)=C(R,t)$。
式~\eqref{eq:moist} 右端的 $D(C)$ 位于散度之内，展开后含
$D'(C)(\partial C/\partial r)^2$ 项，若简单写作 $D\Delta C$ 将丢失该非线性效应。

问题 1 全部物性参数由附录 2 给定（表~\ref{tab:params}），无待定参数，量级与中药材
热风干燥以往数据一致\cite{wang2024drying}。烘房环境由附件 1 给出
（$0\sim14400\,\mathrm{s}$ 共 241 条、间隔 $60\,\mathrm{s}$），本问仅用
$0\sim1800\,\mathrm{s}$ 段，相邻观测点间线性插值，不平滑、不外推；表 1/表 2 所要求
7 个时刻处的烘房条件如图~\ref{fig:chamber} 圆点所示。

\begin{table}[htbp]
  \centering
  \caption{问题 1 物性参数（附录 2）}
  \label{tab:params}
  \begin{tabular}{lll}
    \hline
    参数 & 数值 & 单位 \\
    \hline
    密度 $\rho$ & 820 & $\mathrm{kg/m^3}$ \\
    比热容 $c_p$ & 2600 & $\mathrm{J/(kg\cdot K)}$ \\
    导热系数 $k$ & 0.36 & $\mathrm{W/(m\cdot K)}$ \\
    对流换热系数 $h$ & 25 & $\mathrm{W/(m^2\cdot K)}$ \\
    对流传质系数 $h_m$ & $8\times10^{-7}$ & $\mathrm{m/s}$ \\
    水分扩散系数 $D$ & $7\times10^{-9}\exp(-0.89/C)$ & $\mathrm{m^2/s}$ \\
    半径 $R$／长度 $L$ & 0.02／0.25 & $\mathrm{m}$ \\
    初始温度／水分浓度 & 28／2.55 & $^{\circ}\mathrm{C}$／$\mathrm{kg/kg}$ \\
    \hline
  \end{tabular}
\end{table}

\begin{figure}[htbp]
  \centering
  \includegraphics[width=0.5\textwidth]{fig01_chamber.png}
  \caption{附件 1 烘房条件（阴影为 30 min 考察窗，圆点为表 1/表 2 要求时刻）}
  \label{fig:chamber}
\end{figure}

\subsubsection{模型求解}

本文采用单元中心有限体积法与全隐式时间推进求解式~\eqref{eq:heat}、
\eqref{eq:moist}，具体分为以下四步。

\textbf{(1) 网格。}在归一化坐标 $\xi=r/R$ 下对两方程离散。初始时刻
$C_0\neq C_a(0)$ 使表面出现扩散边界层，厚度按自相似估计为
$\delta_\theta(t)=\sqrt{\pi D t}$（$t=10\,\mathrm{s}$ 时为 $3.9\times10^{-4}\,\mathrm{m}$，
$t=1800\,\mathrm{s}$ 时为 $1.7\times10^{-3}\,\mathrm{m}$）。为分辨该边界层，取表面加密
网格 $\xi_{j+1/2}=1-(1-j/N)^2$（$j=0,1,\dots,N$，$N=800$），最外单元中心距表面
$1.25\times10^{-5}\,\mathrm{m}$，约为最薄边界层的 $1/30$；体积权重采用精确环形体积
$V_i=(\xi_{i+1/2}^2-\xi_{i-1/2}^2)/2$。

\textbf{(2) 界面通量。}界面 $f$ 通量取 $\Phi_f=G_f(u_{f-1}-u_f)$，
$G_f=\xi_f\bar D_f/\Delta\xi$。对 $D=Ae^{-a/C}$，界面系数的精确形式为对数平均
$\bar D_{\log}=(D_1-D_2)/\ln(D_1/D_2)$，与调和平均
$\bar D_{\mathrm{harm}}=2D_1D_2/(D_1+D_2)$ 按 $\varepsilon=\ln(D_2/D_1)$ 展开至
$\varepsilon^2$ 相同，故采用调和平均；Kirchhoff 变换积分通量作为独立校核
\cite{srikiatden2007moisture}。表面按半单元导热与对流传热串联的精确热阻闭性
$G_N=1/[\ln(1/\xi_{c,N})/k+1/(hR)]$，水分边界同理；表面值由通量一致性反解
$T_s=T_a+\Phi_N/(hR)$、$C_s=C_a+\Phi_{C,N}/(h_mR)$。

\textbf{(3) 时间推进。}采用全隐式后向 Euler 与阻尼 Newton 迭代（三对角线性化、解析
Jacobian），以整步与两个半步解之差控制误差并自适应调整步长。温度与水分各取误差预算
$2.5\times10^{-5}$，最大步长 $1\,\mathrm{s}$，在整秒输出点与输入分段节点处截断步长；
两场解耦、分别推进。

\textbf{(4) 输出。}内部点由单元中心值局部二次插值获得，表面点按 (2) 反解，$t=0$ 行
单独保留；结果四舍五入至 4 位小数，按附件 3 模板自 $t=1\,\mathrm{s}$ 起将
$1800\times21$ 时空分布写入 result1.xlsx。

\subsubsection{结果与分析}

按题面格式整理求解结果，如表~\ref{tab:T1}、表~\ref{tab:C1} 所示。

\begin{table}[htbp]
  \centering
  \caption{表 1\quad 30 分钟内药材的温度（$^{\circ}\mathrm{C}$）}
  \label{tab:T1}
  \begin{tabular}{cccccc}
    \hline
    时间/s & \multicolumn{5}{c}{到药材中心的距离/cm} \\
    & 0 & 0.5 & 1 & 1.5 & 2 \\
    \hline
    100   & 28.0001 & 28.0003 & 28.0040 & 28.0327 & 28.1801 \\
    300   & 28.0408 & 28.0635 & 28.1514 & 28.3681 & 28.8489 \\
    600   & 28.4533 & 28.5360 & 28.8039 & 29.3159 & 30.1652 \\
    900   & 29.3243 & 29.4583 & 29.8755 & 30.6161 & 31.7304 \\
    1200  & 30.5427 & 30.7098 & 31.2224 & 32.1126 & 33.4276 \\
    1500  & 31.9957 & 32.1867 & 32.7660 & 33.7463 & 35.1203 \\
    1800  & 33.5753 & 33.7720 & 34.3642 & 35.3621 & 36.7856 \\
    \hline
  \end{tabular}
\end{table}

\begin{table}[htbp]
  \centering
  \caption{表 2\quad 30 分钟内药材的水分浓度（$\mathrm{kg/kg}$）}
  \label{tab:C1}
  \begin{tabular}{cccccc}
    \hline
    时间/s & \multicolumn{5}{c}{到药材中心的距离/cm} \\
    & 0 & 0.5 & 1 & 1.5 & 2 \\
    \hline
    100   & 2.5500 & 2.5500 & 2.5500 & 2.5500 & 2.2470 \\
    300   & 2.5500 & 2.5500 & 2.5500 & 2.5492 & 2.0508 \\
    600   & 2.5500 & 2.5500 & 2.5500 & 2.5353 & 1.8770 \\
    900   & 2.5500 & 2.5500 & 2.5497 & 2.5045 & 1.7547 \\
    1200  & 2.5500 & 2.5500 & 2.5482 & 2.4646 & 1.6586 \\
    1500  & 2.5500 & 2.5499 & 2.5445 & 2.4206 & 1.5787 \\
    1800  & 2.5500 & 2.5497 & 2.5383 & 2.3755 & 1.5102 \\
    \hline
  \end{tabular}
\end{table}

温度场与水分浓度场的时空演化如图~\ref{fig:fields} 所示，表面与中心的时间历程如图
~\ref{fig:series} 所示。

\begin{figure}[htbp]
  \centering
  \includegraphics[width=0.47\textwidth]{fig02_temp_map.png}\hfill
  \includegraphics[width=0.47\textwidth]{fig03_moist_map.png}
  \caption{30 min 内温度场（左）与水分浓度场（右）的时空演化}
  \label{fig:fields}
\end{figure}

\begin{figure}[htbp]
  \centering
  \includegraphics[width=0.5\textwidth]{fig04_timeseries.png}
  \caption{表面与中心温度、水分浓度的时间历程（圆点为表 1/表 2 要求时刻）}
  \label{fig:series}
\end{figure}

\textbf{(1) 温度场。}30 min 内表面温度由 $28\,^{\circ}\mathrm{C}$ 升至
$36.7856\,^{\circ}\mathrm{C}$（升高 $8.79\,^{\circ}\mathrm{C}$），中心升至
$33.5753\,^{\circ}\mathrm{C}$（升高 $5.58\,^{\circ}\mathrm{C}$），$1800\,\mathrm{s}$
时径向温差为 $3.21\,^{\circ}\mathrm{C}$，与 $Bi=1.39$ 的预估一致，集总参数法不适用；
表面温度仍低于烘房 $4.73\,^{\circ}\mathrm{C}$，表明预热平衡在 30 min 内尚未结束，
与 $\tau_T\approx39.5\ \mathrm{min}$ 一致。

\textbf{(2) 水分场。}表面水分浓度由 $2.55\ \mathrm{kg/kg}$ 降至
$1.5102\ \mathrm{kg/kg}$（下降 $40.8\%$），$r=1.5\ \mathrm{cm}$ 处下降 $6.8\%$，
$r=1.0\ \mathrm{cm}$ 处下降 $0.46\%$，中心在 4 位小数精度下不变；显著下降
（$>1\%$）限于表面 $0.5\ \mathrm{cm}$ 薄层，与 $\tau_C\approx22.5\ \mathrm{h}\gg
1800\ \mathrm{s}$ 及 $Bi_m=3.24$ 的预估一致，呈热风干燥中典型的表层先干现象
\cite{wang2024drying,nurhaslina2022review}。表面失水速率随时间衰减：最初 100 s 内
下降 $11.9\%$，$1500\sim1800\,\mathrm{s}$ 内仅下降 $4.3\%$，属对流传质边界驱动的
指数趋近行为。

\subsubsection{模型检验}

为保证表 1/表 2 与 result1.xlsx 中 4 位小数结果的可靠性，本文从以下四方面对模型进行
检验。

\textbf{(1) 温度场解析模态解验证。}问题 1 物性为常数，对分段线性 $T_a(t)$ 存在解析
模态解。令 $\lambda_mJ_1(\lambda_m)=Bi\,J_0(\lambda_m)$、
$A_m=2J_1(\lambda_m)/[\lambda_m(J_0^2(\lambda_m)+J_1^2(\lambda_m))]$，则
\begin{equation}
  T(r,t)=T_a(t)-\sum_{m=1}^{\infty}A_mJ_0\Big(\frac{\lambda_mr}{R}\Big)z_m(t),
  \qquad
  z_m'=T_a'-\frac{\lambda_m^2\alpha}{R^2}z_m,
  \label{eq:modal}
\end{equation}
其中 $z_m$ 在每个 60 s 线性段内逐段精确积分。该解析解与数值解在全部 $1801\times21$
个输出点上比对，最大绝对差为 $5.2916\times10^{-6}\,^{\circ}\mathrm{C}$
（$t=780\ \mathrm{s}$、$r=0.4\ \mathrm{cm}$）；模态截断（180 项与 360 项）差异为
$4.5\times10^{-8}\,^{\circ}\mathrm{C}$。

\textbf{(2) 守恒与物理界。}离散储存量变化与边界交换量的闭合残差为：显热相对残差
$2.0\times10^{-14}$、水分 $5.8\times10^{-13}$；全场温度 $\in[28.0000,36.7856]$、
水分浓度 $\in[1.5102,2.5500]$，无负值、无 NaN。

\textbf{(3) 收敛性与误差预算。}网格 $N$ 逐级加倍，水分场空间收敛阶为 $1.9998$；时间
误差由整步与两个半步解比对控制。误差预算收紧至 $2.5\times10^{-5}$ 后，全场数值误差
估计为温度 $5.336\times10^{-6}\,^{\circ}\mathrm{C}$、水分
$8.105\times10^{-6}\ \mathrm{kg/kg}$，均小于 4 位小数舍入阈值 $5\times10^{-5}$，故
表 1/表 2 与 result1.xlsx 的 4 位小数可靠。上述为给定模型与输入下的数值误差估计，
非实验预测误差。

\textbf{(4) 短时表面自相似核对。}半无限体对流边界的短时自相似解为
$C_0-C_s(t)=(C_0-C_a)\cdot2h_m\sqrt{t/(\pi D)}$\cite{srikiatden2007moisture}：
$t=1\ \mathrm{s}$ 时解析值 $C_s=2.5175$、数值解 $2.5177$，相对差 $8\times10^{-5}$，
表明初始边界层被网格与时间步长（最小时步 $1.0\times10^{-10}\ \mathrm{s}$）充分分辨。

\begin{figure}[H]
  \centering
  \includegraphics[width=0.47\textwidth]{fig06_convergence.png}\hfill
  \includegraphics[width=0.47\textwidth]{fig07_conservation.png}
  \caption{网格收敛性（左）与水分质量守恒（右）}
  \label{fig:verify}
\end{figure}

综上，问题 1 的模型与求解结果已通过解析模态解、守恒与收敛性三重检验，表 1/表 2 与
result1.xlsx 的 4 位小数结果可靠，为问题 2 向完整烘干过程的推广奠定了基础。
